\documentclass[12pt]{article}
\usepackage[T1]{fontenc}
\usepackage[utf8]{inputenc}
\usepackage{lmodern}
\usepackage{textcomp}
\usepackage{enumitem}
\usepackage{geometry}
\usepackage{graphicx}
\usepackage{amsmath, amssymb}
\geometry{margin=1in}
\begin{document}
\begin{center}
Economics - Graduate Level Assessment 20 \
Total Marks: 40 \
Answer all questions.
\end{center}

\section*{Multiple Choice (10 marks)}
\begin{enumerate}[label=\arabic*.]
\item What is the bill rate for a bill with 91 days to maturity and a price of \$ 97,987? Calculate also the coupon-issue yield equivalent for the same bill.
\begin{enumerate}[label=(\alph*)]
    \item Bill rate is 7.70\%, Coupon issue yield is 8.50\%
    \item Bill rate is 8.50\%, Coupon issue yield is 8.24\%
    \item Bill rate is 7.96\%, Coupon issue yield is 8.24\%
    \item Bill rate is 8.24\%, Coupon issue yield is 8.50\%
    \item Bill rate is 8.24\%, Coupon issue yield is 7.96\%
\end{enumerate}
\item Which of the following is a characteristic of monopolistic competition?
\begin{enumerate}[label=(\alph*)]
    \item P = MC.
    \item Perfect competition.
    \item Efficiency.
    \item Inelastic demand.
    \item P =MR.
\end{enumerate}
\item With the presence of a positive externality, which of the following would correct the externality?
\begin{enumerate}[label=(\alph*)]
    \item A government regulation.
    \item An increase in supply.
    \item A government tax.
    \item A government fine.
    \item A decrease in demand.
\end{enumerate}
\item Which of the following is not correct for the perfectly competitive firm, in the long run?
\begin{enumerate}[label=(\alph*)]
    \item price = marginal revenue.
    \item price = minimum average cost.
    \item price = minimum average variable cost.
    \item price = marginal cost.
\end{enumerate}
\item An American buys an entertainment system that was manufactured in China. How does the U.S. national income accounts treat this transaction?
\begin{enumerate}[label=(\alph*)]
    \item Net exports and GDP both fall.
    \item GDP falls and there is no change in net exports.
    \item Net exports rise and there is no change in GDP.
    \item Net exports and GDP both rise.
    \item Net exports rise and GDP falls.
\end{enumerate}
\end{enumerate}

\section*{True / False (10 marks)}
\textit{Answer True or False}
\begin{enumerate}[label=\arabic*.]
\setcounter{enumi}{5}
\item True or False: A trade deficit of \$603 billion suggests that the United States exported more goods than it imported.
\item True or False: Economic resources primarily consist of money, machines, and management rather than traditional factors of production.
\item True or False: The primary language spoken in a country has a direct impact on its export and import volumes.
\item True or False: A growth rate of 10 percent per year would result in the standard of living doubling in 12 years.
\item True or False: Marginal Fixed Cost (MFC) is zero, while Marginal Variable Cost (MVC) is equal to Marginal Total Cost (MC).
\end{enumerate}

\section*{Long Form Response (20 marks)}
\textit{Answer each question with a clear and complete explanation.}
\begin{enumerate}[label=\arabic*.]
\setcounter{enumi}{10}
\item Discuss how positive externalities associated with the production or consumption of a good contribute to market failure, emphasizing the relationship between marginal social benefit and marginal social cost.
\item Analyze the impact of an excise tax imposed on the production of good X in a competitive market, focusing on how the tax burden shifts between consumers and producers based on demand elasticity.
\end{enumerate}

\vfill
\noindent\textit{End of Paper}
\end{document}